\section{Generace video výstupu DVI}

Pro generování DVI signálu je využívána knihovna \texttt{PicoDVI}, která zajišťuje přesné časování,
synchronizaci a samotné generování TMDS signálů potřebných pro přenos obrazu na HDMI.

Výsledný obraz je emulován v rozlišení 320\texttimes240 pixelů. Pro reprezentaci barev byl zvolen
formát RGB565, kde každý pixel zabírá v paměti přesně 16 bitů (2 byty). Tento formát byl
upřednostněn před standardním formátem RGB888 (24 bitů na pixel) z důvodu úspory paměti RAM a
propustnosti sběrnice. Vzhledem k tomu, že původní barevná paleta konzole NES je poměrně omezená,
redukce barevné hloubky nepřináší žádnou znatelnou ztrátu vizuální kvality.

DVI je inicializován ve standardním režimu VGA, který fyzicky generuje rozlišení 640\texttimes480
pixelů při obnovovací frekvenci 60~Hz (viz výpis \ref{lst:dvi_init}). Dle standardu je tento režim
označován jako \textit{Low Pixel Format} s taktovací frekvencí pixelových hodin 25,175~MHz, což
představuje zátěžově nejméně náročný režim podporovaný moderními monitory\cite{dvi-spec}. Vzhledem
k tomu, že interní rozlišení NESu je pouze 256\texttimes240 pixelů, je tento režim pro potřeby
emulátoru dostačující.

\begin{listing}[htbp]
\begin{minted}{c}
#include <dvi.h>
#include <dvi_serialiser.h>

static struct dvi_inst g_dvi;

void core1_main() {
    /* Knihovna PicoDVI přebírá kontrolu nad druhým jádrem
     * procesoru (Core 1), které je dedikováno výhradně pro
     * přesné časování a generování DVI signálu. */
    dvi_register_irqs_this_core(&g_dvi, DMA_IRQ_0);
    dvi_start(&g_dvi);
    dvi_scanbuf_main_16bpp(&g_dvi);
    /* dvi_scanbuf_main_16bpp() běží v nekonečné smyčce
     * a zajišťuje generování obrazu */
    __builtin_unreachable();
}

void init_dvi(void) {
    /* Konfigurace VGA režimu (640x480 pixelů @ 60 Hz) */
    g_dvi.timing = &dvi_timing_640x480p_60hz;
    g_dvi.ser_cfg = DVI_DEFAULT_SERIAL_CONFIG;
    /* Registrace callback funkce pro přípravu nového řádku obrazu */
    g_dvi.scanline_callback = prepare_dvi_scanline;

    dvi_init(
        &g_dvi, next_striped_spin_lock_num(), next_striped_spin_lock_num()
    );
    /* Registrace callback funkce pro přípravu nového řádku obrazu */
    multicore_launch_core1(core1_main);
}
\end{minted}
\caption{Inicializace DVI.}
\label{lst:dvi_init}
\end{listing}

Pro stabilitu video obrazu se udržuje framebuffer s rozlišením 256x240, kde každý pixel je
reprezentován indexem do určité palety. Tento přístup je více úsporný z hlediska využití pamětí, než
udržovat kompletní framebuffer s již dekódovanými pixely.

Knihovna PicoDVI dovoluje nastavení callback funkce, která je automaticky volána vždy, když je
videovýstup připraven na odeslání dalšího řádku. Funkce \texttt{prepare\_dvi\_scanline()} načte z
bufferu odpovídající řádek bytů, pomocí LUT tabulky převede indexy na finální 16bitové barvy a
výsledek předá do fronty řádků PicoDVI. Tento přístup šetří RAM paměť, jelikož skutečné 16bitové
barvy se konstruují dynamicky pouze pro jeden právě vykreslovaný řádek.

Jelikož fyzické DVI rozlišení má šířku 320 pixelů, ale obraz generovaný systémem NES je široký pouze
256 pixelů, je nutné provést horizontální centrování. Toho se dosahuje vložením symetrických černých
pruhů z levé i pravé strany (tzv. \textit{pillarbox}), což zabraňuje nechtěnému roztažení obrazu a
zachovává původní poměr stran (\textit{aspect ratio}). V našem případě pro šetření procesorového
času při dekódování pixelů jsme využili toho, že statické proměnné v C jsou vynulovány, a dekódované
pixely vykreslujeme jen do vyhrazené oblastí v bufferu řádku, proto není potřeba opakovaně centrovat.

Jak již bylo zmíněno v podkapitole \ref{ppu_scanline_callback}, PPU signalizuje dokončení vykreslení
jednoho řádku pomocí definovaného \textit{callbacku}. Zaregistrovaná obslužná funkce
\texttt{prepare\_ppu\_scanline()} ihned překopíruje data z interního bufferu PPU přímo do
příslušného řádku hlavního framebufferu.
